%% =======================================================
%% Actividad 33: "Ventajas y desventajas de la IA en la sociedad"
%% =======================================================

%% ---- CONFIGURACIÓN ------------------------

\documentclass[11pt,letterpaper]{article}

\usepackage[indice]{tecnm}
\usepackage{float}
\usepackage{booktabs}
\usepackage[T1]{fontenc}
\usepackage{xcolor}
\usepackage[most]{tcolorbox} % Carga tcolorbox con soporte para listings y saltos de página
\tcbuselibrary{listings, skins, breakable}

% Paleta de colores (estilo GitHub Light / VS Code)
\definecolor{codebg}{HTML}{F8F9FA}       % Fondo del código
\definecolor{codehead}{HTML}{E9ECEF}     % Fondo de la barra de título
\definecolor{codeborder}{HTML}{CED4DA}   % Borde de la tarjeta
\definecolor{codetitle}{HTML}{212529}    % Texto del título
\definecolor{codekw}{HTML}{0550AE}       % Palabras clave (import, def, for...)
\definecolor{codefunc}{HTML}{8250DF}     % Funciones y métodos (describe, groupby...)
\definecolor{codestr}{HTML}{0A3069}      % Strings
\definecolor{codecomm}{HTML}{6E7781}     % Comentarios
\definecolor{codenum}{HTML}{8C959F}      % Números de línea
\definecolor{codeid}{HTML}{24292F}       % Texto general

% Configuración base de listings
\lstdefinestyle{pythonestilo}{
  language=Python,
  basicstyle=\ttfamily\footnotesize\color{codeid},
  keywordstyle=\color{codekw}\bfseries,
  keywordstyle=[2]\color{codefunc}, % Estilo secundario para métodos de datos
  stringstyle=\color{teal!80!black},
  commentstyle=\color{codecomm}\itshape,
  morekeywords={as, True, False, None, with, yield, lambda},
  morekeywords=[2]{read_csv, head, tail, info, describe, groupby, mean, value_counts, to_csv, print, len, range},
  numbers=left,
  numberstyle=\tiny\ttfamily\color{codenum},
  stepnumber=1,
  numbersep=10pt,
  showstringspaces=false,
  morestring=[b]{"""},
  tabsize=4,
  breaklines=true,
  breakatwhitespace=true,
  columns=flexible,
  keepspaces=true,
  literate=
    {á}{{\'a}}1 {é}{{\'e}}1 {í}{{\'i}}1 {ó}{{\'o}}1 {ú}{{\'u}}1
    {Á}{{\'A}}1 {É}{{\'E}}1 {Í}{{\'I}}1 {Ó}{{\'O}}1 {Ú}{{\'U}}1
    {ü}{{\char252}}1 {«}{{\guillemotleft}}1 {»}{{\guillemotright}}1 {·}{{\textperiodcentered}}1 {ñ}{{\~n}}1 {→}{{$\rightarrow$}}1 {Ñ}{{\~N}}1 {¿}{{?`}}1 {¡}{{!`}}1
}

% Definición del entorno personalizado \begin{codigo}{Título}{etiqueta}
% \newtcblisting[auto counter, number within=section]{codigo}[2][]{
\newtcblisting[auto counter]{codigo}[2][]{
  listing only,
  breakable,                    % Permite que el bloque se divida si pasa de página
  enhanced,
  colback=codebg,
  colframe=codeborder,
  colbacktitle=codehead,
  coltitle=codetitle,
  fonttitle=\bfseries\small\sffamily,
  boxrule=0.6pt,
  arc=4pt,                      % Bordes redondeados
  left=20pt,                    % Espacio izquierdo para los números de línea
  right=8pt,
  top=6pt,
  bottom=6pt,
  toptitle=4pt,
  bottomtitle=4pt,
  before skip=12pt,
  after skip=12pt,
  listing options={style=pythonestilo},
  title={Código \thetcbcounter: #2},
  #1
}

% Paleta para salida de consola (estilo terminal limpia)
\definecolor{termbg}{HTML}{212529}      % Fondo oscuro tipo terminal
\definecolor{termhead}{HTML}{343A40}    % Barra superior
\definecolor{termtext}{HTML}{F8F9FA}    % Texto claro de alta legibilidad
\definecolor{termborder}{HTML}{495057}  % Borde

% Entorno \begin{consola}[label=out:iris]{Título}
\newtcblisting[auto counter]{consola}[2][]{
  listing only,
  breakable,
  colback=termbg,
  colframe=termborder,
  colbacktitle=termhead,
  coltitle=termtext,
  colupper=termtext,
  fonttitle=\bfseries\small\sffamily,
  boxrule=0.6pt,
  arc=4pt,
  left=8pt,
  right=8pt,
  top=6pt,
  bottom=6pt,
  toptitle=4pt,
  bottomtitle=4pt,
  before skip=12pt,
  after skip=12pt,
  listing options={
    basicstyle=\ttfamily\scriptsize\color{termtext}, % \scriptsize evita que las tablas de Pandas se partan de línea
    breaklines=true,
    columns=fixed,        % Mantiene alineadas las columnas de Pandas
    keepspaces=true,
    showstringspaces=false,
    literate=
      {á}{{\'a}}1 {é}{{\'e}}1 {í}{{\'i}}1 {ó}{{\'o}}1 {ú}{{\'u}}1
      {Á}{{\'A}}1 {É}{{\'E}}1 {Í}{{\'I}}1 {Ó}{{\'O}}1 {Ú}{{\'U}}1
      {¿}{{?`}}1 {¡}{{!`}}1 {ü}{{\char252}}1 {«}{{\guillemotleft}}1 {»}{{\guillemotright}}1 {·}{{\textperiodcentered}}1 {ñ}{{\~n}}1 {→}{{$\rightarrow$}}1 {Ñ}{{\~N}}1
  },
  title={Salida \thetcbcounter: #2},
  #1
}

%% ---- DATOS DE LA MATERIA ------------------------

\materia{Inteligencia Artificial y su Ética}
\actividad{Tema 6.2. IA en la sociedad \\ Actividad 33 - Ventajas y desventajas \\ de la IA en la sociedad}
\profesor{M. Lilia Ureña Lugo}


\usepackage{etoolbox}
% ltablex convierte tabularx en longtable, que avanza el contador de tablas una vez de más
\AtEndEnvironment{tabularx}{\addtocounter{table}{-1}}
\newlength{\anchocol}


\begin{document}
\portada
\indicegeneral

\section{Introducción}

La inteligencia artificial ya no se limita a los videojuegos o a la ciencia ficción: hoy está integrada de forma generalizada en la vida cotidiana, muchas veces sin que se note, como parte de sistemas tecnológicos complejos \parencite{iae_u6_t62}. El tema 6.2 está presente en transporte, marketing, salud, finanzas, seguridad y defensa, educación, gestión administrativa, entretenimiento, agricultura y manufactura \parencite{iae_u6_t62}.

Este documento responde a la Actividad 1 del tema 6.2: identificar ventajas y desventajas del uso de la IA en la sociedad (se piden al menos dos de cada una; presento tres de cada una). Todo se apoya en el material del tema \parencite{iae_u6_t62}; no uso otras fuentes. Los libros que el tema cita no los leí, solo recojo lo que el tema dice de ellos.

\section{Desarrollo}

\subsection{Ventajas}

\begin{enumerate}
    \item \textbf{Mejores diagnósticos en salud.} El reconocimiento de imágenes ha permitido avances significativos en el diagnóstico de enfermedades como el cáncer y el Alzheimer \parencite{iae_u6_t62}. Un diagnóstico más preciso o más temprano beneficia directamente a las personas atendidas.
    \item \textbf{Servicios más eficientes y más productividad.} El tema menciona la optimización de servicios públicos y el impulso de aplicaciones comerciales, y señala que el desarrollo tecnológico genera nuevos empleos, impulsa la productividad y promueve el crecimiento económico \parencite{iae_u6_t62}. En la práctica, los algoritmos asumen tareas que antes requerían intervención humana.
    \item \textbf{Educación personalizada.} Con el análisis de grandes volúmenes de datos, la IA permite diseñar estrategias pedagógicas ajustadas a las necesidades individuales de cada estudiante y facilita buscar y adecuar materiales didácticos \parencite{iae_u6_t62}.
\end{enumerate}

\subsection{Desventajas}

\begin{enumerate}
    \item \textbf{Desigualdad y desempleo.} A diferencia de la Revolución Industrial, en la que las máquinas complementaban el trabajo humano, la IA empieza a reemplazarlo, con efectos desproporcionados en ciertos sectores de la población \parencite{iae_u6_t62}. La fuerza laboral tiene que adaptarse y adquirir nuevas habilidades; quien no pueda hacerlo queda en desventaja.
    \item \textbf{Sesgos y discriminación que se refuerzan.} Los datos de entrenamiento reflejan estructuras sociales con sesgos históricos, por lo que aun con algoritmos matemáticamente sólidos esos sesgos se perpetúan y afectan a comunidades vulnerables. Además, las decisiones basadas en datos sesgados pueden generar ciclos de retroalimentación que refuerzan la discriminación, y muchos algoritmos no son de dominio público \parencite{iae_u6_t62}.
    \item \textbf{Privacidad y libertad de expresión.} Personalizar el aprendizaje exige recopilar y analizar grandes volúmenes de datos personales de los estudiantes, lo que genera inquietudes sobre su seguridad y protección. En los medios, el uso de algoritmos para eliminar contenidos sin una evaluación reflexiva preocupa por la libertad de expresión \parencite{iae_u6_t62}. En política, los sesgos y la proliferación de bots pueden comprometer la diversidad y la equidad del debate público \parencite{iae_u6_t62}.
\end{enumerate}

\begin{table}[H]
\centering
\scriptsize
\renewcommand{\arraystretch}{1.25}
\begin{tabularx}{\textwidth}{@{}p{3cm}XX@{}}
\toprule
\textbf{Ámbito} & \textbf{Ventaja} & \textbf{Desventaja} \\
\midrule
Salud & Diagnóstico de cáncer y Alzheimer por reconocimiento de imágenes. & Datos de salud muy delicados; el sesgo puede afectar a comunidades vulnerables. \\
Educación & Aprendizaje personalizado y búsqueda ágil de materiales. & Recopilación masiva de datos de estudiantes y riesgo para su privacidad. \\
Trabajo y economía & Productividad, crecimiento y nuevos empleos. & Reemplazo de funciones, desempleo y desigualdad. \\
Medios y política & Información acorde con los intereses del usuario; detección de noticias falsas. & Eliminación automática de contenidos, bots y amplificación de sesgos. \\
\bottomrule
\end{tabularx}
\caption{Ventajas y desventajas por ámbito, según el tema 6.2 \parencite{iae_u6_t62}.}\label{tab:ambitos}
\end{table}

\section{Conclusión}

Las ventajas y las desventajas aparecen en los mismos ámbitos: la personalización que mejora la enseñanza es la misma que exige datos personales, y la automatización que da productividad es la que reemplaza empleos. Por eso, más que preguntar si la IA es buena o mala, el tema propone preguntar quién se beneficia y quién decide su rumbo \parencite{iae_u6_t62}. Mi conclusión es que sus ventajas son reales, pero que su reparto depende de decisiones humanas: quién aplica modelos más equitativos, con qué propósito y bajo qué reglas. El tema lo resume diciendo que ser conscientes de estas implicaciones permite decidir hasta dónde queremos que la IA cambie nuestra forma de vida \parencite{iae_u6_t62}.

\referencias

\end{document}
